\chapter{如何展示，而不只是告知}

這場講座幾乎沒有通常意義上的數學，卻仍有一條形式性的脊骨。它要我們把散文想成一套推理系統。我們究竟是直接把結論交給讀者，還是把足以讓讀者自行回收結論的證據交給讀者？整場講座鋪排得相當細密：一段刻意寫得不好的開場文字先製造問題，一連串讓步又防止我們把這堂課變成一句口號，接著六條原則逐步磨利技藝。我們會依照這個順序來走，因為整場論證不只依賴個別主張，也同樣依賴它自身的推進方式。

\begin{remark}
這裡沒有真正的黑板數學可供轉錄。下列方程式都是編者加上的示意圖式：它們把講座的邏輯壓縮起來，但並不意圖暗示畫面上曾出現任何符號形式化。
\end{remark}

\section{開場的失敗}

講座不是從定義開始，而是從一場失敗開始。我們讀到一段荒野場景：一個女孩被說成很害怕，四周的聲音被說成令人恐懼，而地面則被說成像一位守護者。所有重要的事情都被宣布了，但真正變得可感的，卻極少。

這是正確的教學開場。講者想讓我們先感受到它的無力，然後才告訴我們它叫什麼。那段文字說出了恐懼，卻沒有把恐懼戲劇化；它說出了氛圍，卻沒有把氛圍造出來。我們得到的是情感判決，而不是足以支撐那些判決的感官壓力與敘事壓力。

\subsection{Question \& Answer}

講座在這裡停下來，提出了支配全章的那個問題。

這幅圖景究竟哪裡出了問題？

問題不在於這段文字用了情緒詞。真正的麻煩在於：它還沒提供證據，就先把結論端出來了。恐懼是被宣告的，而不是被演出的；安全感也是被宣告的，而不是被掙得的。就在場景本應使那些感受變得無可避免的時刻，讀者卻被直接告知該怎麼想、怎麼感覺。

\begin{definition}
就本講座的意義而言，\emph{告知} 是直接陳述結論、概括、特質或情緒；而 \emph{展示} 則是提供戲劇性的、感官性的、行為性的或視角性的證據，使同樣的結論能夠被推知出來。
\end{definition}

K.\ M.\ Weiland 那個簡潔的區分，替講座立下第一條清楚的坐標軸：
\begin{align}
\text{展示} &= \text{戲劇化},\\
\text{告知} &= \text{概括化}.
\end{align}

這其實已經修正了一個常見混淆。概括本身不是敵人；沒有支撐的概括才是敵人。接下來的講座，將花其餘的時間來告訴我們：什麼時候概括在做有用的工作，什麼時候它又太早地把場景封死了。

\section{展示、告知，以及讀者所做的工作}

在把問題尖銳地提出之後，講座立刻放鬆了規則。告知本身並不壞。若我們試圖展示每一個稍縱即逝的事實、每一段流逝的時間、每一個例行動作，小說就會膨脹、鬆垮、失去形狀。故事需要壓縮，也需要橋樑；它既需要大筆掃過，也需要真正被居住的時刻。

這也就是為什麼《The Secret Garden》的例子出現在這裡。文本會告訴我們，Mary Lennox 是個看起來惹人厭的孩子，而且常常生病；但它並不止於這個判詞，而是立刻轉向她瘦削的臉、瘦削的身體、淡淡的頭髮，以及發酸的表情。概括被證據錨住了。這正是講座心目中正當混合的模型。

講座的中心邏輯，可以示意地寫成
\begin{equation}
\text{具體細節} \leadsto \text{讀者推知} \leadsto \text{情緒或判斷}.
\end{equation}

也正是在這裡，Andrew Stanton 的說故事比喻開始變得有用。逐字稿在這一段略有雜訊，但核心意思非常清楚。我們若想要觀眾積極參與，就不應直接把 ``4'' 交給他們，而應該給他們 ``2+2''，讓他們自己做最後那一步：
\begin{equation}
2+2 \to \text{觀眾推知}, \qquad 4 \to \text{過早封閉}.
\end{equation}

講座堅持，這不是裝飾性的建議。讀者其實願意工作，只要那份工作被巧妙地藏在經驗裡。他們想要推理、補全、發現。從這個意義上說，展示是一種有組織的不完整：我們故意拿掉最後那個標籤，讓讀者在場景的壓力之下，自己走到那裡。

\begin{figure}[h]
\centering
\begin{tikzpicture}[>=stealth]
\node[draw, rectangle, rounded corners, minimum width=3.4cm, minimum height=1.1cm, align=center] (a) at (0,0) {具體細節\\或證據};
\node[draw, rectangle, rounded corners, minimum width=3.2cm, minimum height=1.1cm, align=center] (b) at (4.7,0) {讀者\\推知};
\node[draw, rectangle, rounded corners, minimum width=3.6cm, minimum height=1.1cm, align=center] (c) at (9.5,0) {情緒、判斷\\或氛圍};
\draw[->, thick] (a) -- (b);
\draw[->, thick] (b) -- (c);
\end{tikzpicture}
\caption{講座中的基本推理管線。我們不是先命名判決，而是先給出證據，使判決幾乎不可避免地浮現。}
\end{figure}

從這裡開始，講座把主張擴展開來。對電影敘事成立的事，對散文也同樣成立。當意義是被暗示出來的，而不是被手指著指出來的時候，讀者會更在意。也正因此，講座現在才能退一步，追問這句格言究竟從何而來。

\section{為什麼會有這句格言}

在進入六條原則之前，講座先插入一段短短的歷史繞行。這不是裝飾，而是在解釋：為什麼這條規則會有這麼大的力量。

講者把現代這句口號放進寫實主義的傳統裡：那是一種試圖以足夠的即時性與誠實，來描寫普通生活的努力，使它不再顯得被浪漫化，或被修辭過度操控。接著，Percy Lubbock 被引入作為這種區分的重要形式化者。他的關鍵想法是：當故事被當成某種應該被展示的東西來處理，以至於它彷彿在自己講述自己時，小說才真正成為藝術。

Janice Hardy 那種實用性的重述，則把這段文學史重新翻回工作坊語言。只要一句話聽起來像是角色的思想、知覺或壓力，我們大致就是安全的；但一旦它聽起來像作者跳進來替場景解碼，沉浸感便會削弱。

我們可以把這種實用傾向概括為
\begin{equation}
\text{作者性解釋} \downarrow \;\Longrightarrow\; \text{讀者沉浸} \uparrow.
\end{equation}

這不是一條絕對律；講座稍後還會把說明與敘述放回它們正當的位置。但這段繞行把真正的目標澄清了。問題不在於告知這件事本身存在，而在於：就在我們想讓讀者從內部活過場景的那一刻，作者卻變得太可見了。

有了這層較大的動機，講座便已準備好走向實作。它現在依序轉向六條指導原則。

\section{證據、具體細節與感官特異性}

前三條原則形成了一個自然的群組。它們從不同角度回答的是同一個問題：如果我們不直接宣告情感性或描述性的結論，那麼要拿什麼來取代它？

\subsection{先給證據，再給判詞}

第一步最乾淨。當敘述者說某個丈夫心地善良，當主角相信一位朋友有罪，或者某一句寫著 Adam knew Gwen liked him，講者要我們在這個結論前暫停一下，去尋找它的根據。究竟是什麼讓任何人這樣想？

講座裡引用的那個 thought-verb 練習，說的也是同一件事。為了修改，我們暫時不信任像 \emph{knew}、\emph{believed}、\emph{wanted}、\emph{realized} 與 \emph{imagined} 這樣的動詞。不是因為它們被禁止，而是因為它們常常把真正的場景藏了起來。

\paragraph{Worked revision.} 講座中的例句是 ``Adam knew Gwen liked him.'' 修補的方法並不神祕，而是按部就班的。
\begin{enumerate}
\item 先從結論開始：Adam knew Gwen liked him.
\item 接著追問 Adam 實際上感知到了什麼。
\item 用一次次在置物櫃邊的相遇、漆面金屬上的鞋跟印、香水味，以及仍帶著溫度的密碼鎖，取代那個藏起來的判斷。
\item 到讀者能在無須幫助的情況下自行抵達 Adam 的結論時，就停下來。
\end{enumerate}

若示意地寫成
\begin{align}
\text{``Adam knew Gwen liked him''}
&\xrightarrow{\text{修訂}}
\text{置物櫃} + \text{鞋跟印} + \text{香水} + \text{溫熱的鎖}.
\end{align}

所以，第一條原則並不只是 ``多加細節'' 而已，而是更精確地說：用允許該結論成立的證據，來取代那個結論本身。

\subsection{從抽象走向具體}

第二條原則，則把同一套邏輯應用到情緒詞與評價性形容詞上。Harvey Chapman 的例子在這裡非常有效。與其說 Toby 走回家時比他此生任何時候都更快樂，不如給我們那個怎麼也下不來的傻笑，和他走路時那種彈跳、幾乎要絆倒的輕盈。

這可以寫成一條一般性的修訂規則：
\begin{equation}
\text{抽象標籤} \xrightarrow{\text{修訂}} \text{可觀察的證據}.
\end{equation}

講座裡幾個替換例子，清楚展示了這個模式：
\begin{align}
\text{快樂}
&\xrightarrow{\text{展示}}
\text{傻氣的笑容} + \text{彈跳般的步伐},\\
\text{陰森}
&\xrightarrow{\text{展示}}
\text{一片森林低鳴著早已死去孩童的哭聲}.
\end{align}

重點並不是說抽象總是不正當，而是說：抽象字往往來得太早。它們在場景尚未真正讓意義變得可感時，就搶先告訴我們這個場景「代表什麼」。講座中那個陰森森林的例子之所以有用，也正因如此。``The dark forest felt eerie'' 並不是假的，只是太弱。修訂過的句子並沒有反駁那個形容詞；它只是把它掙了出來。

\subsection{Question \& Answer}

我們怎麼知道自己仍然在抽象裡寫作？

講座給出的操作性答案很簡單：問問看，攝影機看不看得見它。這當然不是完整檢驗，因為散文還可以喚起嗅覺、觸覺、味覺與聲音。但它是很強的第一道過濾。若攝影機看不見，那我們就該追問：文本究竟提供了什麼感官證據？

這便直接引出第三條示意規則：
\begin{equation}
\text{原因陳述} \xrightarrow{\text{展示}} \text{可見或可感的效果}.
\end{equation}

Jerry Jenkins 的例子在這裡幾乎是理想範例：
\begin{align}
\text{寒冷}
&\xrightarrow{\text{展示}}
\text{豎起的衣領} + \text{拉緊的圍巾} + \text{插進口袋的手} + \text{迎風轉開的臉},\\
\text{失明}
&\xrightarrow{\text{展示}}
\text{白手杖探尋長椅},\\
\text{深秋}
&\xrightarrow{\text{展示}}
\text{腳下踩碎的落葉聲}.
\end{align}

這是講座中最重要的一條局部推導之一。我們不再命名原因，而是寫下它留下的痕跡。不說 ``it was cold''，而是展示寒冷會讓人做什麼；不說 ``it was late fall''，而是讓地面替我們發出那個季節的聲音。

講座接著把這個要點延伸到可見性之外。Delilah Dawson 那場市場戲之所以更有力，不只是因為它具體，而是因為它變得在感官上與地方上都具體：肉桂、咖啡、絲穗、粉末狀的番紅花。當場景的細節既鮮活，又只屬於它自己的世界時，沉浸感就會變強。

最後，講者又補上一個細緻的警告。即使一個場景在技術上是可視的，它仍可能毫無生氣，若其中細節過於可預期。一場下雨的葬禮可能是清楚的，卻仍然陳腐。這也就是為什麼講座特地停在 Gail Carson Levine 的建議上：若某個場景讓人覺得老套，就改變場景、改變情緒角度，或引入一個打破期待的局部異常。到了這裡，我們不再只是用細節取代抽象，而是在用屬於 \emph{這個} 故事的細節，取代通用性的細節。

\section{超越制式身體語言}

到了這裡，講座做出一個重要修正。一旦作者被告知要展示，他們往往會抓住最現成的速記方式：緊握的拳頭、交叉的雙臂、咬緊的牙關、狂跳的心臟、翻騰的胃。這些並非毫無用處，但也極其容易被濫用。

此時的危險，比先前更精確了。先前的危險是赤裸裸的抽象；現在的危險則是一種淺層的具體：它確實給出了一個可見的徵象，暗示某種情緒正在發生，卻不足以告訴我們那究竟是 \emph{哪一種} 情緒。

\subsection{Question \& Answer}

光靠身體語言就夠了嗎？

通常不夠。它或許能指示 \emph{有什麼} 正在發生，卻常常無法傳達 \emph{究竟是哪一種} 發生。若只把憤怒、挫折、羞辱、嫉妒、驚恐或悲傷全都壓縮成一個制式手勢，它們看起來就會彼此相似。講座的重點是：身體語言本身往往太過壓縮，無法單獨承載情緒上的細微差別。

Robin Patchen 那個前後對照的例子，幾乎把這一點講死。第一版裡，Mary 看見時鐘，心臟幾乎要從胸口跳出來，胃裡翻騰，她害怕出了什麼事。乍看之下，這好像已經在展示了，因為我們確實看見了身體反應。但講者堅持，這場景仍然顯得誇張而過度解釋。身體在做全部工作，思想過程卻幾乎不見。

修訂版則把重心移開。陽光流進窗內。Mary 在休息充足中醒來。她看見時鐘，發現小 Jane 居然一整夜都睡著了。緊接著，另一段記憶闖入：Billy。然後行動幾乎立刻跟上：她掀開被子，抓起睡袍，在心裡重複醫生的保證。現在，情緒開始有了輪廓。這不再只是警慌，而是處在記憶與否認壓力之下的警慌。

講座把這壓縮成一條值得保留的因果鏈：
\begin{equation}
\text{思想} \to \text{情緒} \to \text{行動}.
\end{equation}

這也許是整場講座最乾淨的一條 ``方程式''。它告訴我們，當一個場景的情緒顯得薄弱時，應該去哪裡找。若我們太快跳向身體速記，就會略過那個讓情緒獲得精確形狀的思想；一旦把思想補回來，接下來的行動也就變得更具體、更可信。

因此，一條有用的修訂序列可以是：
\begin{enumerate}
\item 用局部的知覺、辨認與記憶序列，取代泛泛的內臟式驚慌。
\item 讓場景本身參與進來：開著的窗、陽光、時鐘、睡袍。
\item 讓動詞承擔力道：\emph{flipped} 與 \emph{snatched} 都比中性的動作動詞更有力量。
\item 也讓句子結構本身幫忙；短單位與突兀轉折，都可以直接演出迫切感。
\end{enumerate}

講座接著把這轉化成實際建議。當場景帶有高度情緒時，我們往往更應該放慢它。讓我們聽見角色在想什麼，讓幾個細節停留，讓行動從中冒出來。只要思想與行動真的在場，讀者是聰明到足以自己還原情緒的。

\section{對話、預告式寫法與冗餘}

第五條原則把同樣的紀律轉移到言說上。對話本身就可以直接展示情緒，只要我們足夠信任那些話，以及整個場景，能夠自己完成工作。

講座先從一個明顯的例子講起。若 Mary 正在對 Bob 生氣，那麼那句話本身就可以承載怒氣。一旦台詞本身已經完成了那份情緒，像 ``angrily'' 這樣的標記，往往只會削弱而不會幫忙。它告訴了我們台詞早就已經展示出來的事。

\subsection{Question \& Answer}

對話什麼時候開始變成預告，而不再是展示？

當敘述提前宣布一個動機、意圖或語氣結論，而這些東西本來就已經存在於對話交換本身之中時，對話就開始變成預告式寫法。若我們先寫他 tried to be diplomatic，然後又給他一行和緩的話，我們就把資訊加倍了。若我們寫她 changed the subject，接著又立刻讓她把話題轉走，我們就在解釋讀者本可以自行推知的事。

若示意地寫成，
\begin{align}
\text{對話} + \text{冗餘的解釋}
&\to \text{預告式寫法},\\
\text{對話} + \text{精確的動作拍點}
&\to \text{更銳利的弦外之音}.
\end{align}

講座中的修訂版對話，正遵守了這條規則。它不再用說明性旁白，而是給出一個經過選擇的動作拍點與一個經過選擇的標記：他捏了捏鼻樑；她低聲說。台詞本身仍承載著主要力量，而周圍的細節則把社會性的幾何關係磨得更清楚，同時又不替我們解碼。

講座接著再往前推一步。若我們想讓情緒性的對話更有力，把場景幾乎當成一場戲劇來起草，往往會有幫助。這個限制會迫使我們把語氣直接寫在交換本身裡。Oscar Wilde 在這裡被提起，正因為他的對話能做到這一點。簡短的反駁、難以置信的重複、一個小小的措辭轉折：這些都能讓一種情緒先變得可聽見，然後敘述才來得及解釋。

這也為最後一條原則準備了過渡。若我們能更信任對話，那麼我們也能更廣泛地信任視角。冗餘的問題，會在說明與描寫裡再次回來。

\section{敘事聲音、說明與節奏}

第六條原則，把講座從局部修補擴展到敘事建築。展示不只是把一個形容詞換成一個圖像而已，它同時也是視角、說明與節奏的問題。

講座先從幾組成對例子開始。``He was rude and inconsiderate'' 變成一行對著一位推嬰兒車、正掙扎著的女人吼出來的話。``She was uncomfortable around him'' 變成 ``She stiffened in his embrace.'' ``The house was huge'' 變成一個大到足夠容納整個家族的廚房。``She was hungry'' 則變成一碗幾乎被吸進去的湯。這些句子之所以更強，不只是因為它們具體，而是因為它們帶著視角。每一句都是從一顆被放進特定情境中的心智抵達的。

同樣的原則也支配著說明。Janice Hardy 對「對讀者有益」與「對角色有益」的區分，在這裡至關重要。那些只是為了向讀者簡報而存在的資訊，常會聽起來像是從場景外部塞進來的；但如果資訊是經由角色實際會注意、會害怕、會發問或會說出口的方式被過濾出來，它就仍留在故事的經驗場域之內。

那個餐廳例子，把這個洞見磨得更尖。雨、鬆餅、一個信封、滴答作響的時鐘：外部事實本身未必需要改變多少。可是，一旦視角改變，整個場景便會改變。對一顆軍事化的心智而言，雨聲會變成槍火；對一個受驚的女孩而言，窗面則變成模糊與封閉。這不是修飾性的差別，而是敘事聲音在做結構性工作。

同樣的壓力也解釋了為什麼 ``as you know Bob'' 式對話總會失敗。當一個角色向另一個角色解釋雙方都已經知道的事，那句話服務的就不是角色，而是太過明目張膽地在服務讀者。作者於是再度變得可見。

到了這裡，講座做出最後、也是最重要的一次修正。真正的教訓不是 ``永遠展示''，而是 ``展示，而不只是告知''。告知仍然是我們的節奏工具之一。它讓我們能夠搭橋、壓縮、概括與前進。展示則讓我們能夠停留、居住與加壓。

我們可以把這個平衡示意地寫成
\begin{equation}
\text{告知} + \text{具體細節} = \text{當節奏需要時，一種強而有力的混合}.
\end{equation}

對講座節奏規則做進一步壓縮，則可以寫成
\begin{align}
\text{場景重要性} \uparrow &\Rightarrow \text{展示密度} \uparrow,\\
\text{橋接或壓縮需求} \uparrow &\Rightarrow \text{告知比例} \uparrow.
\end{align}

講座最後那張 ``cheat sheet'' 可以概括如下。

\begin{table}[h]
\centering
\begin{tabular}{|p{0.42\textwidth}|p{0.42\textwidth}|}
\hline
\textbf{需要速度時，用告知或混合} & \textbf{需要臨場感時，用展示} \\
\hline
從 A 點到 B 點的橋接 & 情緒 \\
\hline
例行事務、時間流逝、重複性的對話 & 感官：視覺、聽覺、嗅覺、味覺、觸覺 \\
\hline
大範圍說明或偶爾的背景交代 & 對人或場所的假設與判斷 \\
\hline
某些世界觀建構式概述 & 對話的語氣、動機與交談意圖 \\
\hline
大尺度的思考或敘事掃掠 & 那些讓情緒得以被推知的具體痕跡 \\
\hline
\end{tabular}
\caption{講座最後的節奏平衡。告知並沒有被廢除，而是被分配。}
\end{table}

這也就是為什麼講座可以毫不矛盾地在結尾讚美：在說明性開頭裡，強而有力的敘事聲音依舊能成立。只要那個聲音本身帶著張力與壓力，一個以敘述開始的段落也仍可能非常有力。告知並沒有被逐出小說，它只是被規訓、被安置。

這個最後的平衡，也替我們準備好重返開頭那段荒野文字。Delia Owens 的那段文字之所以成立，是因為它不只是宣布恐懼消退、安慰抵達。它讓土地在女孩跌倒時接住她，讓痛楚像水滲進沙裡一樣消散，也讓沼澤成為母親。情感轉變不再是從外部被報告，而是透過圖像、運動與聲音被演出來。講座在它開始的地方結束，但那個壞例子，現在已經被轉化成一條標準。

\section{Summary}

現在，我們可以用較成熟的形式來重述講座的論證。展示不是拒絕告知；展示是一種紀律：在場景本身已經提供了足以讓判決浮現的證據之前，不先把最終判詞交出去。

當一行文字顯得薄弱時，講座給了我們一串嚴格的追問：
\begin{itemize}
\item 我們是不是在提供證據之前，就先提供了結論？
\item 我們是不是命名了一種情緒，而在其實只需一個具體痕跡時就搶先下了標籤？
\item 攝影機看得見我們所寫的東西嗎？或者感官至少能否註冊它？
\item 我們是不是倚賴了制式身體語言，而不是思想、場景與行動？
\item 我們是不是在替本來已經會自己說明的對話做解說？
\item 說明是在服務角色視角，還是在單純替讀者做簡報？
\end{itemize}

因此，最後的原則既平衡又實際。我們在需要臨場感、壓力與讀者參與時展示；在需要壓縮、轉場或大尺度敘事掃掠時告知。技藝的所在，就是知道哪些時刻值得給出 ``2+2''，而哪些時刻則可以誠實地直接給成 ``4''。